% ══════════════════════════════════════════════════════════════
% Chapter 23: Kubernetes Deployment
% ══════════════════════════════════════════════════════════════
\chapter{Deployment: Kubernetes}
\label{ch:k8s-deployment}

\begin{learnbox}
\begin{itemize}
  \item Deploy the blockchain network on Kubernetes for production-grade orchestration
  \item Configure StatefulSets with per-pod storage for stateful blockchain nodes
  \item Understand service discovery, DNS, and networking in Kubernetes
  \item Implement horizontal pod autoscaling and high availability
  \item Migrate from Docker Compose to Kubernetes manifests
  \item Troubleshoot and operate Kubernetes deployments
\end{itemize}
\end{learnbox}

\lettrine{D}{o}\index{Docker}cker Compose excels for local development, but \index{Kubernetes}Kubernetes is built for long-running systems in production. If a \index{blockchain}blockchain node crashes, \index{Kubernetes}Kubernetes restarts it automatically. If you need more nodes, you change a replica count. If you push a new image, \index{Kubernetes}Kubernetes rolls it out without downtime. This chapter shows how the same \index{Docker}Docker images from Chapter~\ref{ch:docker-deployment} are deployed into a system that can run reliably in production.

\begin{infobox}{Prerequisites}
This chapter assumes you have read Chapter~\ref{ch:docker-deployment} (Docker Compose) and have basic familiarity with Kubernetes concepts. A local cluster like Minikube or kind is recommended for following along.
\end{infobox}

% ──────────────────────────────────────────────────────────────
\section{Introduction and Architecture}
\label{sec:k8s-introduction}

\subsection{Why Kubernetes After Docker Compose}

\begin{table}[htbp]
\caption{Docker Compose vs. Kubernetes}
\label{tab:compose-vs-k8s}
\centering\small
\begin{tabularx}{0.92\textwidth}{lll}
\toprule
\textbf{Feature} & \textbf{Docker Compose} & \textbf{Kubernetes} \\
\midrule
Self-healing & Limited & \checkmark Restarts and reschedules \\
Rolling updates & Manual & \checkmark Built-in rollout/rollback \\
Service discovery & Docker DNS & \checkmark Cluster DNS + Services \\
Autoscaling & Manual & \checkmark HPA \\
Stateful workloads & Possible & \checkmark First-class StatefulSets \\
Multi-node & No & \checkmark Native \\
\bottomrule
\end{tabularx}
\end{table}

In this repo, \index{Kubernetes}Kubernetes matters because \textbf{\index{mining}miners and webservers are stateful} (disk-backed databases and \index{wallet}wallets), and \index{Kubernetes}Kubernetes provides StatefulSets with per-pod PersistentVolumeClaims---the right primitive to manage this cleanly.

\subsection{What You'll Build}

\begin{figure}[htbp]
  \centering
  \begin{tikzpicture}[
    node distance = 1.0cm and 0.9cm,
    cluster/.style = {rectangle, rounded corners=8pt, draw=sectioncolor, fill=sectioncolor!5,
                      minimum width=6.5cm, minimum height=7cm},
    section/.style = {rectangle, rounded corners=5pt, draw=sectioncolor!60, fill=codebg,
                      text width=2.8cm, align=center, minimum height=0.7cm, font=\small\sffamily},
    item/.style = {rectangle, rounded corners=3pt, draw=sectioncolor!40, fill=white,
                   text width=2.0cm, align=center, minimum height=0.5cm, font=\tiny\ttfamily},
  ]

  % Cluster border
  \node[cluster] (cluster) at (3.25, 3.5) {};
  \node[font=\small\sffamily\bfseries, anchor=north west] at (0.5,6.8) {Kubernetes Cluster (namespace: blockchain)};

  % StatefulSet: Miner
  \node[section] (miner_ss) at (1.2, 5.5) {
    \textbf{StatefulSet:} \\
    miner
  };
  \node[item] (miner0) at (0.4, 4.6) {miner-0};
  \node[item] (miner1) at (1.2, 4.6) {miner-1};
  \node[item] (pvc_m) at (2.0, 4.6) {PVCs};

  % StatefulSet: Webserver
  \node[section] (web_ss) at (5.3, 5.5) {
    \textbf{StatefulSet:} \\
    webserver
  };
  \node[item] (web0) at (4.5, 4.6) {webserver-0};
  \node[item] (web1) at (5.3, 4.6) {webserver-1};
  \node[item] (pvc_w) at (6.1, 4.6) {PVCs};

  % Deployment: Redis
  \node[section] (redis_deploy) at (3.25, 3.8) {
    \textbf{Deployment:} \\
    redis
  };

  % Services
  \node[section] (services) at (1.2, 2.8) {
    \textbf{Services}
  };
  \node[item] (svc_m_hl) at (0.2, 2.0) {miner-hl};
  \node[item] (svc_m) at (1.2, 2.0) {miner-svc};
  \node[item] (svc_w_hl) at (2.2, 2.0) {web-hl};

  % ConfigMaps
  \node[section] (cm) at (5.3, 2.8) {
    \textbf{ConfigMaps}
  };
  \node[item] (cm_cfg) at (4.5, 2.0) {config};
  \node[item] (cm_rate) at (6.1, 2.0) {rate-limit};

  % HPA
  \node[section] (hpa) at (3.25, 1.0) {
    \textbf{HPA:} miner-hpa, \\
    webserver-hpa
  };

  % Secrets
  \node[section] (secrets) at (3.25, 0.1) {
    \textbf{Secrets:} \\
    blockchain-secrets
  };

  \end{tikzpicture}
  \caption{Kubernetes Deployment Architecture: StatefulSets, Services, and Configuration}
  \label{fig:k8s-architecture}
\end{figure}

% ──────────────────────────────────────────────────────────────
\section{Core Kubernetes Concepts}
\label{sec:k8s-concepts}

\subsection{Namespace}

A Namespace provides resource isolation and scoping. All resources in this repo live in the \code{blockchain} namespace.

\begin{listing}[htbp]
\caption{Namespace Definition}
\label{lst:ch23-namespace}
\begin{minted}[fontsize=\footnotesize]{yaml}
apiVersion: v1
kind: Namespace
metadata:
  name: blockchain
  labels:
    name: blockchain
    environment: production
\end{minted}
\end{listing}

\subsection{ConfigMap}

ConfigMaps store non-sensitive configuration data. In this repo, two ConfigMaps are used:
\begin{itemize}
  \item \code{blockchain-config}: General node settings (as env vars)
  \item \code{rate-limit-settings}: \code{Settings.toml} (mounted as file)
\end{itemize}

\begin{listing}[htbp]
\caption{ConfigMap for Node Settings}
\label{lst:ch23-configmap}
\begin{minted}[fontsize=\footnotesize]{yaml}
apiVersion: v1
kind: ConfigMap
metadata:
  name: blockchain-config
  namespace: blockchain
data:
  MINER_NODE_IS_MINER: "yes"
  MINER_NODE_CONNECT_NODES: "local"
  WEBSERVER_NODE_IS_WEB_SERVER: "yes"
  WEBSERVER_NODE_CONNECT_NODES: >-
    miner-service.blockchain.svc.cluster.local:2001
  SEQUENTIAL_STARTUP: "no"
  WALLET_FILE: "wallets/wallets.dat"
\end{minted}
\end{listing}

\subsection{Secret}

Secrets store sensitive data like API keys. Values are base64-encoded (not encrypted by default).

\begin{listing}[htbp]
\caption{Secret for API Keys}
\label{lst:ch23-secret}
\begin{minted}[fontsize=\footnotesize]{yaml}
apiVersion: v1
kind: Secret
metadata:
  name: blockchain-secrets
  namespace: blockchain
type: Opaque
stringData:
  BITCOIN_API_ADMIN_KEY: "admin-secret"
  BITCOIN_API_WALLET_KEY: "wallet-secret"
  # Optional: if omitted, entrypoint auto-creates and persists
  # MINER_ADDRESS: ""
\end{minted}
\end{listing}

\subsection{StatefulSet vs. Deployment}

A StatefulSet manages \textbf{stateful pods} with:
\begin{itemize}
  \item Stable pod identity (\code{miner-0}, \code{miner-1}, \ldots)
  \item Stable per-pod storage via \code{volumeClaimTemplates}
  \item Stable per-pod DNS identity when used with a headless Service
\end{itemize}

A Deployment manages \textbf{stateless pods} where replicas are interchangeable. In this repo, both miners and webservers are stateful (disk-backed chain DB and wallets), so both run as StatefulSets.

\begin{listing}[htbp]
\caption{StatefulSet for Webservers}
\label{lst:ch23-statefulset}
\begin{minted}[fontsize=\footnotesize]{yaml}
apiVersion: apps/v1
kind: StatefulSet
metadata:
  name: webserver
  namespace: blockchain
spec:
  serviceName: webserver-headless
  replicas: 2
  selector:
    matchLabels:
      app: webserver
  template:
    metadata:
      labels:
        app: webserver
    spec:
      containers:
      - name: blockchain-node
        image: blockchain-node:latest
        ports:
        - name: web
          containerPort: 8080
        - name: p2p
          containerPort: 2001
        volumeMounts:
        - name: webserver-data
          mountPath: /app/data
        - name: webserver-wallets
          mountPath: /app/wallets
  volumeClaimTemplates:
  - metadata:
      name: webserver-data
    spec:
      accessModes: ["ReadWriteOnce"]
      resources:
        requests:
          storage: 50Gi
  - metadata:
      name: webserver-wallets
    spec:
      accessModes: ["ReadWriteOnce"]
      resources:
        requests:
          storage: 1Gi
\end{minted}
\end{listing}

The \code{volumeClaimTemplates} section creates \textbf{one PVC per pod}, ensuring each webserver/miner has isolated disk storage. This avoids shared-disk locking issues (e.g., Sled DB lock contention) that would occur if multiple replicas shared the same storage path.

\subsection{Service: Normal and Headless}

In Kubernetes, it is common to have \textbf{two Services} for the same application:

\begin{enumerate}
  \item A \textbf{headless Service} (\code{clusterIP: None}) for stable \textbf{per-pod identity} (DNS)
  \item A \textbf{normal Service} (ClusterIP/NodePort/LoadBalancer) for stable \textbf{load-balanced access}
\end{enumerate}

\begin{listing}[htbp]
\caption{Headless Service (Per-Pod DNS Identity)}
\label{lst:ch23-service-headless}
\begin{minted}[fontsize=\footnotesize]{yaml}
apiVersion: v1
kind: Service
metadata:
  name: webserver-headless
  namespace: blockchain
spec:
  clusterIP: None  # Headless service
  selector:
    app: webserver
  ports:
  - name: web
    port: 8080
    targetPort: 8080
  - name: p2p
    port: 2001
    targetPort: 2001
\end{minted}
\end{listing}

A headless Service enables stable DNS names like:
\begin{equation}
\text{\code{webserver-0.webserver-headless.blockchain.svc.cluster.local}}
\end{equation}

\begin{listing}[htbp]
\caption{Normal Service (Load-Balanced Access)}
\label{lst:ch23-service-normal}
\begin{minted}[fontsize=\footnotesize]{yaml}
apiVersion: v1
kind: Service
metadata:
  name: webserver-service
  namespace: blockchain
spec:
  type: LoadBalancer
  selector:
    app: webserver
  ports:
  - name: web
    port: 8080
    targetPort: 8080
\end{minted}
\end{listing}

\subsection{Service Discovery}

In Kubernetes, pods connect using Service DNS names:
\begin{equation}
\text{\code{<service-name>.<namespace>.svc.cluster.local}}
\end{equation}

\textbf{Example}: A webserver connecting to a miner:

\begin{listing}[H]
\caption{Service Discovery: Webserver to Miner}
\label{lst:ch23-service-discovery}
\begin{minted}[fontsize=\footnotesize]{bash}
# Docker Compose
NODE_CONNECT_NODES=miner_1:2001

# Kubernetes
NODE_CONNECT_NODES=miner-service.blockchain.svc.cluster.local:2001
\end{minted}
\end{listing}

% ──────────────────────────────────────────────────────────────
\section{Setup and Quick Start}
\label{sec:k8s-quickstart}

\subsection{Prerequisites}

Install the required tools:

\begin{listing}[htbp]
\caption{Install Kubernetes Tools (macOS)}
\label{lst:ch23-install-tools}
\begin{minted}[fontsize=\footnotesize]{bash}
brew install kubectl
brew install minikube
brew install docker
\end{minted}
\end{listing}

\subsection{Start Minikube Cluster}

\begin{listing}[H]
\caption{Start Local Kubernetes Cluster}
\label{lst:ch23-start-minikube}
\begin{minted}[fontsize=\footnotesize]{bash}
cd ci/kubernetes

# Start cluster with metrics-server addon (required for HPA)
minikube start --cpus=4 --memory=3072mb \
  --addons=metrics-server
\end{minted}
\end{listing}

\subsection{Build and Load Docker Image}

\begin{listing}[H]
\caption{Build Image into Minikube}
\label{lst:ch23-build-image}
\begin{minted}[fontsize=\footnotesize]{bash}
# Use Minikube's Docker daemon
eval $(minikube docker-env)

# Build from repository root
cd ../../  # (where Dockerfile is)
docker build -t blockchain-node:latest -f Dockerfile .

# Restore normal Docker
eval $(minikube docker-env -u)
\end{minted}
\end{listing}

\subsection{Deploy Using Script}

\begin{listing}[H]
\caption{Deploy All Kubernetes Resources}
\label{lst:ch23-deploy}
\begin{minted}[fontsize=\footnotesize]{bash}
cd ci/kubernetes/manifests
./deploy.sh
\end{minted}
\end{listing}

The deployment script applies manifests in dependency order:
\begin{enumerate}
  \item Namespace
  \item ConfigMaps and Secrets
  \item PersistentVolumeClaims
  \item Redis (rate limiting backend)
  \item StatefulSets (miners and webservers)
  \item Services
  \item HPA configurations
  \item Pod Disruption Budgets
\end{enumerate}

\subsection{Verify Deployment}

\begin{listing}[H]
\caption{Verify Pods Are Running}
\label{lst:ch23-verify}
\begin{minted}[fontsize=\footnotesize]{bash}
# Check pod status
kubectl get pods -n blockchain

# Wait for all pods to be ready
kubectl wait --for=condition=ready pod \
  -l app=miner -n blockchain --timeout=300s
kubectl wait --for=condition=ready pod \
  -l app=webserver -n blockchain --timeout=300s

# View all resources
kubectl get all -n blockchain
\end{minted}
\end{listing}

\subsection{Access Webserver (Port Forward)}

\begin{listing}[H]
\caption{Port-Forward to Webserver Service}
\label{lst:ch23-port-forward}
\begin{minted}[fontsize=\footnotesize]{bash}
kubectl port-forward -n blockchain \
  svc/webserver-service 8080:8080

# Then access at http://localhost:8080
\end{minted}
\end{listing}

% ──────────────────────────────────────────────────────────────
\section{Migration from Docker Compose}
\label{sec:k8s-migration}

\subsection{Step 1: Build and Push Image}

For local Minikube: use the image loaded in Step 3 above.

For cloud: tag and push to your registry:

\begin{listing}[htbp]
\caption{Push to Container Registry}
\label{lst:ch23-push-image}
\begin{minted}[fontsize=\footnotesize]{bash}
docker build -t blockchain-node:latest -f Dockerfile .
docker tag blockchain-node:latest your-registry/blockchain-node:v1.0.0
docker push your-registry/blockchain-node:v1.0.0
\end{minted}
\end{listing}

\subsection{Step 2: Create Namespace}

\begin{listing}[H]
\caption{Create Blockchain Namespace}
\label{lst:ch23-create-ns}
\begin{minted}[fontsize=\footnotesize]{bash}
kubectl apply -f 01-namespace.yaml
\end{minted}
\end{listing}

\subsection{Step 3: Create ConfigMaps and Secrets}

\begin{listing}[H]
\caption{Apply Configuration and Secrets}
\label{lst:ch23-config-secrets}
\begin{minted}[fontsize=\footnotesize]{bash}
kubectl apply -f 02-configmap.yaml
kubectl apply -f 03-secrets.yaml
kubectl apply -f 14-configmap-rate-limit.yaml
\end{minted}
\end{listing}

\subsection{Step 4: Create StatefulSets}

StatefulSets automatically create per-pod PVCs based on \code{volumeClaimTemplates}:

\begin{listing}[htbp]
\caption{Deploy StatefulSets}
\label{lst:ch23-deploy-statefulsets}
\begin{minted}[fontsize=\footnotesize]{bash}
kubectl apply -f 15-redis.yaml
kubectl apply -f 06-statefulset-miner.yaml
kubectl apply -f 07-deployment-webserver.yaml
\end{minted}
\end{listing}

\subsection{Step 5: Create Services}

\begin{listing}[H]
\caption{Deploy Services for Networking}
\label{lst:ch23-deploy-services}
\begin{minted}[fontsize=\footnotesize]{bash}
kubectl apply -f 08-service-miner-headless.yaml
kubectl apply -f 08-service-miner.yaml
kubectl apply -f 09-service-webserver-headless.yaml
kubectl apply -f 09-service-webserver.yaml
\end{minted}
\end{listing}

\subsection{Step 6: Configure HPA (Autoscaling)}

\begin{listing}[H]
\caption{Deploy Horizontal Pod Autoscalers}
\label{lst:ch23-deploy-hpa}
\begin{minted}[fontsize=\footnotesize]{bash}
kubectl apply -f 10-hpa-webserver.yaml
kubectl apply -f 11-hpa-miner.yaml
\end{minted}
\end{listing}

\subsection{Key Differences from Docker Compose}

\begin{table}[htbp]
\caption{Migration: Docker Compose to Kubernetes}
\label{tab:migration-differences}
\centering\small
\begin{tabularx}{0.92\textwidth}{lll}
\toprule
\textbf{Aspect} & \textbf{Docker Compose} & \textbf{Kubernetes} \\
\midrule
Service discovery & \code{miner\_1:2001} & \code{miner-service\ldots:2001} \\
Scaling & \code{--scale miner=3} & \code{kubectl scale statefulset miner \ldots} \\
Port mapping & Manual in compose & Service types (ClusterIP, LB) \\
Storage & Docker volumes & PersistentVolumeClaims \\
Configuration & Env vars & ConfigMaps + Secrets \\
Rolling updates & Manual & \code{kubectl rollout} \\
\bottomrule
\end{tabularx}
\end{table}

% ──────────────────────────────────────────────────────────────
\section{Operations and Scaling}
\label{sec:k8s-operations}

\subsection{Manual Scaling}

Scale StatefulSets by changing replica count:

\begin{listing}[htbp]
\caption{Manual Pod Scaling}
\label{lst:ch23-scale}
\begin{minted}[fontsize=\footnotesize]{bash}
# Scale webservers to 5 replicas
kubectl scale statefulset webserver -n blockchain \
  --replicas=5

# Scale miners to 3
kubectl scale statefulset miner -n blockchain \
  --replicas=3
\end{minted}
\end{listing}

\subsection{Viewing Logs}

\begin{listing}[H]
\caption{Stream Container Logs}
\label{lst:ch23-logs}
\begin{minted}[fontsize=\footnotesize]{bash}
# All miner logs
kubectl logs -n blockchain -l app=miner -f

# Specific pod
kubectl logs -n blockchain miner-0 -f

# Show previous container logs (for crashes)
kubectl logs -n blockchain miner-0 --previous
\end{minted}
\end{listing}

\subsection{Update Configuration}

\begin{listing}[H]
\caption{Update and Rollout Configuration}
\label{lst:ch23-update-config}
\begin{minted}[fontsize=\footnotesize]{bash}
# Edit ConfigMap
kubectl edit configmap blockchain-config -n blockchain

# Restart pods to pick up changes
kubectl rollout restart statefulset/webserver \
  -n blockchain
kubectl rollout restart statefulset/miner -n blockchain
\end{minted}
\end{listing}

\subsection{Rolling Updates}

\begin{listing}[H]
\caption{Rolling Update Strategy}
\label{lst:ch23-rolling-update}
\begin{minted}[fontsize=\footnotesize]{bash}
# Update image
kubectl set image statefulset/webserver \
  blockchain-node=blockchain-node:v1.1.0 \
  -n blockchain

# Watch rollout progress
kubectl rollout status statefulset/webserver \
  -n blockchain

# Rollback if needed
kubectl rollout undo statefulset/webserver -n blockchain
\end{minted}
\end{listing}

% ──────────────────────────────────────────────────────────────
\section{Horizontal Pod Autoscaling (HPA)}
\label{sec:k8s-autoscaling}

HPA automatically adjusts replica counts based on CPU and memory metrics:

\begin{listing}[htbp]
\caption{HPA Configuration}
\label{lst:ch23-hpa}
\begin{minted}[fontsize=\footnotesize]{yaml}
apiVersion: autoscaling/v2
kind: HorizontalPodAutoscaler
metadata:
  name: webserver-hpa
  namespace: blockchain
spec:
  scaleTargetRef:
    apiVersion: apps/v1
    kind: StatefulSet
    name: webserver
  minReplicas: 1
  maxReplicas: 10
  metrics:
  - type: Resource
    resource:
      name: cpu
      target:
        type: Utilization
        averageUtilization: 70
\end{minted}
\end{listing}

\subsection{Verify HPA Status}

\begin{listing}[H]
\caption{Check Autoscaling Status}
\label{lst:ch23-hpa-status}
\begin{minted}[fontsize=\footnotesize]{bash}
# View HPA status
kubectl describe hpa webserver-hpa -n blockchain

# Check current resource usage
kubectl top pods -n blockchain

# If \`kubectl top\` fails, enable metrics-server:
minikube addons enable metrics-server

# Verify Metrics API is available
kubectl get apiservices | grep metrics
\end{minted}
\end{listing}

% ──────────────────────────────────────────────────────────────
\section{Troubleshooting}
\label{sec:k8s-troubleshooting}

\subsection{Symptom: CrashLoopBackOff}

Container starts and immediately exits, then Kubernetes restarts it in a loop.

\begin{listing}[htbp]
\caption{Debug CrashLoopBackOff}
\label{lst:ch23-crashloop}
\begin{minted}[fontsize=\footnotesize]{bash}
# View logs from crashed container
kubectl logs -n blockchain miner-0 --previous \
  --tail 200

# Describe pod for events and probe failures
kubectl describe pod miner-0 -n blockchain

# Common cause: invalid MINER\_ADDRESS
# Fix: remove placeholder or set real address
\end{minted}
\end{listing}

\subsection{Symptom: Pod Running but Not Ready}

Container is running, but readiness probes fail.

\begin{listing}[htbp]
\caption{Debug Readiness Issues}
\label{lst:ch23-not-ready}
\begin{minted}[fontsize=\footnotesize]{bash}
# Describe pod to see probe failures
kubectl describe pod miner-0 -n blockchain

# Common causes:
# - Port mismatches between container and probe
# - Process not listening on expected port
# - Application startup delay
\end{minted}
\end{listing}

\subsection{Symptom: Storage Lock Errors}

Multiple webserver pods trying to open same Sled database simultaneously.

\begin{listing}[htbp]
\caption{Debug Storage Lock Errors}
\label{lst:ch23-storage-lock}
\begin{minted}[fontsize=\footnotesize]{bash}
# Verify each webserver has its own PVC
kubectl get pvc -n blockchain | grep webserver

# If sharing storage, ensure StatefulSet with
# volumeClaimTemplates is used (not Deployment)
# with shared PVC
\end{minted}
\end{listing}

\subsection{Symptom: HPA Not Working}

\begin{listing}[H]
\caption{Debug HPA Issues}
\label{lst:ch23-hpa-debug}
\begin{minted}[fontsize=\footnotesize]{bash}
# Check HPA status
kubectl describe hpa webserver-hpa -n blockchain

# Check if metrics are available
kubectl top pods -n blockchain

# If metrics unavailable, enable metrics-server:
minikube addons enable metrics-server

# Verify the Metrics API is registered:
kubectl get apiservices | grep metrics
\end{minted}
\end{listing}

% ──────────────────────────────────────────────────────────────
\section{Cleanup}
\label{sec:k8s-cleanup}

\begin{listing}[H]
\caption{Delete Kubernetes Deployment}
\label{lst:ch23-cleanup}
\begin{minted}[fontsize=\footnotesize]{bash}
cd ci/kubernetes/manifests
./undeploy.sh

# Or delete namespace (removes all resources)
kubectl delete namespace blockchain

# Stop local cluster (keeps cluster on disk)
minikube stop

# Remove local cluster completely
minikube delete
\end{minted}
\end{listing}

% ──────────────────────────────────────────────────────────────
\section{Production Considerations}
\label{sec:k8s-production}

\subsection{Persistent Storage}

In production:
\begin{itemize}
  \item Use cloud-provided storage classes (AWS EBS, Google Persistent Disks, Azure Disks)
  \item Ensure PVCs are backed by reliable block storage with snapshots/backups
  \item StatefulSets with per-pod PVCs prevent shared-disk locking
\end{itemize}

\subsection{Resource Requests and Limits}

Set appropriate CPU and memory requests/limits to:
\begin{itemize}
  \item Enable proper scheduling
  \item Prevent resource starvation
  \item Trigger autoscaling based on actual usage
\end{itemize}

\begin{listing}[htbp]
\caption{Resource Configuration}
\label{lst:ch23-resources}
\begin{minted}[fontsize=\footnotesize]{yaml}
resources:
  requests:
    cpu: "500m"
    memory: "512Mi"
  limits:
    cpu: "2000m"
    memory: "2Gi"
\end{minted}
\end{listing}

\subsection{Network Policies}

In production, restrict network traffic using Kubernetes NetworkPolicies to limit which pods can communicate.

\subsection{Secrets Management}

For production, use proper secret management:
\begin{itemize}
  \item Sealed Secrets (automatic encryption)
  \item External Secrets Operator (integrate with Vault, AWS Secrets Manager)
  \item Avoid base64-encoding as security (not encryption)
\end{itemize}

% ──────────────────────────────────────────────────────────────
\begin{chaptersummary}
\item We deployed the blockchain network on Kubernetes using StatefulSets for stateful miners and webservers.
\item We understood StatefulSets with per-pod PersistentVolumeClaims as the right primitive for disk-backed databases.
\item We implemented service discovery using normal and headless Services for load-balanced and per-pod DNS access.
\item We configured Horizontal Pod Autoscaling to automatically adjust replica counts based on CPU/memory metrics.
\item We migrated from Docker Compose to Kubernetes by building images, creating manifests, and applying resources in dependency order.
\item We scaled deployments manually using \code{kubectl scale} and automatically using HPA.
\item We troubleshot common issues (CrashLoopBackOff, readiness failures, storage locks, HPA misconfiguration).
\item We considered production best practices: persistent storage, resource management, network policies, and secrets management.
\end{chaptersummary}

---


